% ============================================================================
% main.tex -- Exemplo de uso do tema beamerthemeIFCEBaturite
%
% ESTE ARQUIVO É SÓ UMA DEMONSTRAÇÃO DO TEMPLATE: os textos/imagens abaixo
% são conteúdo dummy (lorem ipsum, "fotos" genéricas em TikZ) só para dar
% uma ideia realista do resultado final -- não é uma lista de exercícios
% real. Cada lista futura deve ser um projeto separado, reaproveitando
% este .sty.
%
% A configuração (ativos institucionais + metadados) fica em config.tex,
% não aqui -- este arquivo só tem conteúdo (capa + questões).
%
% Compilar com XeLaTeX (necessário para as fontes Carlito / Open Sans):
%   xelatex main.tex
% ============================================================================
\documentclass{beamer}
\usetheme{IFCEBaturite} % define o tamanho de papel correto (16:9, escala do PPTX)
% ============================================================================
% config.tex -- configuração desta lista (edite isto, não o main.tex)
%
% Tudo que muda de lista pra lista -- ou que é fixo mas fica fora do
% repositório -- vai aqui. O main.tex (ou o .tex renomeado da lista) só
% deve ter conteúdo: capa + questões.
% ============================================================================

% --- Ativos institucionais (logo + ícone, sempre os mesmos) -----------------
% Ficam num arquivo externo, fora do controle de versão -- o mesmo
% caminho serve pra toda lista, então normalmente você NÃO precisa
% mexer nessa linha depois de configurá-la uma vez.
%
% Ajuste o caminho abaixo para o arquivo .tex externo que contém a sua
% chamada de \IFCEativosInstitucionais com os caminhos reais das
% imagens, por exemplo:
%
%   % conteúdo de /home/santiago/algum-lugar/institucional.tex:
%   \IFCEativosInstitucionais
%     {\includegraphics[width=7.25cm]{/home/santiago/algum-lugar/logo-ifce-baturite.pdf}}
%     {\includegraphics[width=0.37cm]{/home/santiago/algum-lugar/icone-ifce-mini.pdf}}
%
% Se o arquivo não existir (como no próprio template, sem nenhuma lista
% real configurada ainda), caem placeholders genéricos no lugar --
% a compilação nunca quebra por causa disso.
\InputIfFileExists{/CAMINHO/EXTERNO/institucional.tex}{}{%
  \IFCEativosInstitucionais
    {\IFCEplaceholder[4.0cm]{logo da instituição}}
    {\textcolor{ifcegreen}{\rule{0.3cm}{0.3cm}}}
}

% --- Metadados desta lista (edite a cada lista nova) ------------------------
% Os valores abaixo são os do exemplo de demonstração deste próprio
% template (main.tex); numa lista real, troque todos os 8 campos.
\IFCEmetadados
  {IFCE Baturité}                          % instituição
  {Técnico Integrado em Comércio}           % curso/turma
  {Física II/Complementar}                  % disciplina
  {Lista de Exercícios}                     % tipo de conteúdo
  {Enem e UECE}                             % descrição do conteúdo
  {Vetores}                                 % assunto/tópico
  {Paulo Santiago}                          % professor
  {\IFCEplaceholder[2.0cm]{logo do curso}}  % logomarca do curso (troque por
                                             % \includegraphics{...} quando tiver a arte)


\begin{document}

% ---------------------------------------------------------------------
% Slide de título
% ---------------------------------------------------------------------
\CapaIFCE{\ImagemCapa}

% ---------------------------------------------------------------------
% Slides de questão -- numeração "01", "02", ... automática
% ---------------------------------------------------------------------
\question
  {ENEM 2023}
  {Um bloco de massa $m$ é abandonado do repouso no topo de um plano
   inclinado sem atrito, que faz um ângulo $\theta$ com a horizontal.
   Sobre o bloco atuam três forças: o peso $\vec{P}$, a força normal
   $\vec{N}$ e, ao longo do movimento, uma força de resistência do ar
   $\vec{F}_r$, sempre oposta à velocidade.

   Considerando a decomposição do peso nas direções paralela e
   perpendicular ao plano, é correto afirmar que a componente que
   acelera o bloco ao longo da rampa é}
  {\IFCEimagemDummy{16.0cm}{10.5cm}}

\question
  {UECE 2022.1}
  {Dois vetores $\vec{A}$ e $\vec{B}$ possuem módulos $8\,\text{N}$ e
   $6\,\text{N}$, respectivamente, e formam entre si um ângulo de
   $90^\circ$. Determine o módulo do vetor resultante
   $\vec{R} = \vec{A} + \vec{B}$ e o ângulo que $\vec{R}$ forma com
   $\vec{A}$.}
  {\textbf{Resolução}

   Como os vetores são perpendiculares, aplica-se o Teorema de
   Pitágoras para obter o módulo da resultante:
   \[
     R = \sqrt{A^2 + B^2} = \sqrt{8^2 + 6^2} = \sqrt{100} = 10\,\text{N}
   \]
   e o ângulo $\alpha$ entre $\vec{R}$ e $\vec{A}$ satisfaz
   $\tan\alpha = B/A = 6/8 = 0{,}75$, logo $\alpha \approx 36{,}9^\circ$.}

\end{document}
